\chapter[Band 1: Grundlagen der Logik]{Band 1 auf einen Blick}
\noindent\textbf{Grundlagen der Logik}\par\medskip
\noindent\emph{Lesart.} \(P,Q,R\) sind Formeln, \(t,u\) Terme und
\(\Gamma\) eine Menge offener Annahmen. Dieser Grundlagenband legt die
Sprache und den Kalkül fest. Rechts stehen daher ausdrücklich auch
\emph{Schlussregeln}; sie sind keine bereits bewiesenen semantischen Sätze.
\begin{BandUebersicht}
\textbf{Terme und atomare Formeln}
\UebFormel{\begin{aligned}
&x,\ c,\ f(t_1,\ldots,t_n)\quad\text{Terme},\\
&P(t_1,\ldots,t_n),\ t=u\quad\text{Formeln}.
\end{aligned}}
Funktions- und Prädikatensymbole haben jeweils eine feste Stelligkeit.
\UebQuellen{Band 1, Abschnitte \emph{Terme} und \emph{Aussagen der Logik}.}
&
\textit{Gleichheitsregeln}
\UebFormel{\begin{aligned}
&\vdash t=t,\\
&t=u\dsep P(t)\vdash P(u).
\end{aligned}}
Bei der Ersetzung werden gebundene Variablen nicht eingefangen.
\UebQuellen{Band 1, Regeln \(=I\) und \(=E\).}\\
\addlinespace[.8ex]
\textbf{Zusammengesetzte Formeln}
\UebFormel{\begin{aligned}
&\neg P,\quad P\land Q,\quad P\lor Q,\\
&P\rightarrow Q,\quad P\leftrightarrow Q.
\end{aligned}}
Diese Bildungsregeln bestimmen die zulässigen Junktorenformeln.
&
\textit{Einführung und Elimination der Konjunktion}
\UebFormel{\begin{aligned}
&P\dsep Q\vdash P\land Q,\\
&P\land Q\vdash P,\qquad P\land Q\vdash Q.
\end{aligned}}
\UebQuellen{Band 1, Regeln \(\land I\), \(\land E_1\), \(\land E_2\).}\\
\addlinespace[.8ex]
\textbf{Beweis und Ableitbarkeit}
Eine endliche Ableitung endet mit der behaupteten Formel; jede Zeile
ist Annahme, Axiom oder durch eine Regel gerechtfertigt.
\UebFormel{\Gamma\vdash_{\mathcal C}P.}
Ein Satz hat keine offenen Annahmen. Inhaltliche Axiome werden später
als Voraussetzungen einer Theorie eingeführt.
&
\textit{Implikation und Entlassung einer Annahme}
\UebFormel{\begin{aligned}
&[P]\vdots Q\quad\vdash\quad P\rightarrow Q,\\
&P\rightarrow Q\dsep P\vdash Q.
\end{aligned}}
Im ersten Schema wird die ausgezeichnete Annahme \(P\) entlassen;
andere offene Annahmen bleiben bestehen.
\UebQuellen{Band 1, Regeln \(\rightarrow I\) und \(\rightarrow E\).}\\
\addlinespace[.8ex]
\textbf{Fallunterscheidung und Widerspruch}
Die Metanotation \([P]\vdots R\) bezeichnet einen Hilfsbeweis mit
ausgezeichneter Annahme \(P\). \(\bot\) bezeichnet Falschheit.
&
\textit{Disjunktions- und Negationsregeln}
\UebFormel{\begin{aligned}
&P\vdash P\lor Q,\qquad Q\vdash P\lor Q,\\
&P\lor Q\dsep[P]\vdots R\dsep[Q]\vdots R\\
&\hspace{18mm}\vdash R,\\
&P\dsep\neg P\vdash\bot,\\
&[P]\vdots\bot\quad\vdash\quad\neg P.
\end{aligned}}
\UebQuellen{Band 1, Regeln \(\lor I\), \(\lor E\), \(\bot I\), \(\neg I\).}\\
\addlinespace[.8ex]
\textbf{Quantifizierte Formeln}
\UebFormel{\forall x\,P(x),\qquad\exists x\,P(x).}
Ein Term muss für die ersetzte Variable einsetzbar sein. Bei den
Eigenvariablenregeln gelten die im Band angegebenen Frischebedingungen.
&
\textit{Quantorenregeln}
\UebFormel{\begin{aligned}
&\forall x\,P(x)\vdash P(t),\\
&P(t)\vdash\exists x\,P(x),\\
&P(x)\mid_x\quad\vdash\quad\forall x\,P(x),\\
&\exists x\,P(x)\dsep[P(x)]\vdots Q\mid_x\\
&\hspace{18mm}\vdash Q.
\end{aligned}}
Bei \(\forall I\) darf \(x\) in keiner offenen Annahme frei sein;
bei \(\exists E\) auch nicht in der Schlussformel \(Q\).
\UebQuellen{Band 1, Regeln \(\forall I\), \(\forall E\), \(\exists I\), \(\exists E\).}\\
\addlinespace[.8ex]
\textbf{Definitionen und eindeutige Beschreibung}
Explizite Definitionen führen Abkürzungen ein; implizite Definitionen
legen eine Struktur durch Axiome fest. Ein Iota-Term setzt eindeutige
Existenz voraus.
\UebFormel{a\coloneqq\iota x\,P(x).}
&
\textit{Verwendung einer eindeutigen Beschreibung}
Unter der Voraussetzung \(\exists!x\,P(x)\) kann der eindeutig
beschriebene Gegenstand eingesetzt werden:
\UebFormel{\begin{aligned}
&P(\iota x\,P(x)),\\
&P(y)\vdash y=\iota x\,P(x).
\end{aligned}}
\UebQuellen{Band 1, Abschnitt \emph{Das Iota-Symbol und Iota-Definitionen}.}\\
\addlinespace[.8ex]
\textbf{Metanotation}
\UebFormel{\begin{aligned}
&a\neq b\quad\text{steht für}\quad\neg(a=b),\\
&P\eqvdash Q\quad\text{steht für}\\
&\qquad P\vdash Q\ \text{und}\ Q\vdash P.
\end{aligned}}
Metasymbole strukturieren die Darstellung von Ableitungen.
&
\textit{Ersetzung äquivalenter Teilausdrücke}
Bewiesene Äquivalenzen erlauben die im Band formulierte
metasprachliche Ersetzung in einem Formelschema. Dabei werden
die Voraussetzungen der verwendeten Äquivalenz mitgeführt.
\UebQuellen{Band 1, Abschnitt \emph{Meta-Regel: Ersetzung äquivalenter Teilausdrücke}.}\\
\end{BandUebersicht}
